\textbf{Theorem (Markov’s inequality).}  
Let \((\Omega,\mathcal{F},\mathbb{P})\) be a probability space and let  
\(X:\Omega\to[0,\infty]\) be a non‑negative random variable.  
For any real number \(a>0\),
\[
\mathbb{P}\bigl(X\ge a\bigr)\;\le\;\frac{\mathbb{E}[X]}{a}.
\]

\begin{proof}
\paragraph{1. Trivial and degenerate cases.}  
If \(X\equiv0\) a.s., then \(\mathbb{E}[X]=0\) and \(\mathbb{P}(X\ge a)=0\) for every
\(a>0\); the inequality holds with equality.  
Henceforth assume \(a>0\) and that \(X\) is an arbitrary non‑negative random
variable (not identically zero).

\paragraph{2. Indicator of the event \(\{X\ge a\}\).}  
Define the indicator function
\[
\mathbf{1}_{\{X\ge a\}}(\omega)=
\begin{cases}
1, & X(\omega)\ge a,\\[2mm]
0, & X(\omega)< a,
\end{cases}
\qquad (\omega\in\Omega).
\]
Since \(\{X\ge a\}\in\mathcal{F}\), this function is measurable.

\paragraph{3. Pointwise bound.}  
For every \(\omega\in\Omega\) we have
\begin{equation}
\mathbf{1}_{\{X\ge a\}}(\omega)\;\le\;\frac{X(\omega)}{a}. \tag{1}
\end{equation}
Indeed, if \(X(\omega)<a\) then the left–hand side is \(0\le X(\omega)/a\);
if \(X(\omega)\ge a\) then \(X(\omega)/a\ge1\), whence \(1\le X(\omega)/a\).

\paragraph{4. Taking expectations.}  
Because (1) holds pointwise, the monotonicity of the Lebesgue integral gives
\[
\mathbb{E}\bigl[\mathbf{1}_{\{X\ge a\}}\bigr]\;\le\;\mathbb{E}\!\left[\frac{X}{a}\right].
\]
Linearity of the expectation allows us to pull out the constant \(1/a\):
\begin{equation}
\mathbb{E}\bigl[\mathbf{1}_{\{X\ge a\}}\bigr]\;\le\;\frac{1}{a}\,\mathbb{E}[X]. \tag{2}
\end{equation}

\paragraph{5. Identification of the left‑hand side.}  
For any event \(A\in\mathcal{F}\),
\[
\mathbb{E}\bigl[\mathbf{1}_{A}\bigr]=\mathbb{P}(A).
\]
With \(A=\{X\ge a\}\) we obtain
\[
\mathbb{E}\bigl[\mathbf{1}_{\{X\ge a\}}\bigr]=\mathbb{P}\bigl(X\ge a\bigr).
\]
Substituting this into (2) yields the desired inequality
\[
\boxed{\;\mathbb{P}\bigl(X\ge a\bigr)\le\frac{\mathbb{E}[X]}{a}\;}.
\]

\paragraph{6. Equality conditions.}  
Equality in (2) can occur only if every inequality used in the proof is an
equality.

\begin{enumerate}
\item \textbf{Equality in the pointwise bound \eqref{1}.}  
    This requires
    \begin{equation}
    \mathbf{1}_{\{X\ge a\}}(\omega)=\frac{X(\omega)}{a}
    \quad\text{for }\mathbb{P}\text{-a.e. }\omega. \tag{3}
    \end{equation}
\item \textbf{Solving \eqref{3}.}  
    If \(X(\omega)<a\), then \eqref{3} forces \(X(\omega)=0\).  
    If \(X(\omega)\ge a\), then \eqref{3} forces \(X(\omega)=a\).
    Hence \eqref{3} holds a.s. iff \(X\) takes only the two values \(0\) and
    \(a\) (with arbitrary probabilities).  Formally,
    \begin{equation}
    \mathbb{P}\bigl(X\in\{0,a\}\bigr)=1. \tag{4}
    \end{equation}
    The degenerate case \(X\equiv0\) is included in (4) by allowing the
    probability of the value \(a\) to be zero.
\end{enumerate}
Thus equality in Markov’s inequality holds exactly when \(X\) is almost surely
supported on \(\{0,a\}\).

\end{proof}